% --- LaTeX CV Template - S. Venkatraman ---

% --- Set document class and font size ---

\documentclass[letterpaper, 11pt]{article}

% --- Package imports ---

\usepackage{hyperref, enumitem, longtable, amsmath, array}
\usepackage{fancyhdr}
\usepackage{ifthen}
\setlength{\headheight}{13.59999pt}

% --- Public/Private toggle ---
\newboolean{privateversion}
\setboolean{privateversion}{true} % Set to 'true' for private version, 'false' for public version

% Define commands to conditionally show content
\newcommand{\ifprivate}[2]{\ifthenelse{\boolean{privateversion}}{#1}{#2}} % Usage: \ifprivate{private content}{public content}
\newcommand{\privateonly}[1]{\ifprivate{#1}{}} % Usage: \privateonly{content that only appears in private version}
\newcommand{\publiconly}[1]{\ifprivate{}{#1}} % Usage: \publiconly{content that only appears in public version}

% --- Page layout settings ---

% Set page margins
\usepackage[left=0.7in, right=0.8in, bottom=.8in, top=0.8in, headsep=0.25in, footskip=.2in]{geometry}

% Set line spacing
\renewcommand{\baselinestretch}{1.1}

% --- Page formatting settings ---

% Set link colors
\usepackage[dvipsnames]{xcolor}
\definecolor{CustomBlue}{HTML}{3e4bb3} % Define your custom color
\hypersetup{colorlinks=true, linkcolor=MidnightBlue, urlcolor=MidnightBlue}

% Set font to Libertine, including math support
\usepackage{libertine}
\usepackage[libertine]{newtxmath}

% Enable page numbering
\pagenumbering{arabic}

% Define header and footer style
\pagestyle{fancy}
\fancyhf{} % Clear all header and footer fields
\fancyhead[L]{Dr. Alexander Nettekoven} % Left header text
\fancyhead[R]{Curriculum Vitae} % Right header text
\fancyfoot[C]{\thepage} % Center footer with page number

% Define first page style (no header)
\fancypagestyle{firstpage}{
  \fancyhf{} % Clear all header and footer fields
  \fancyfoot[C]{\thepage} % Only keep page number in footer
}

% Define font size and color for section headings
\newcommand{\headingfont}{\Large\color{CustomBlue}}

% --- CV section settings ---

% Note: each section of this table (Education, Awards, Publications etc.) is 
% stored in a two-column table. The left-hand column is narrow (1 inch) and is 
% meant to store dates. The right-hand column is wide (5.2 inches) and stores 
% the main text.  Sections in which each entry might have multiple lines 
% (e.g., Education) are stored in a 'SectionTable' environment). Sections in 
% which each entry might just have one line are stored in a 'SectionTableSingleSpace'
% environment. The only difference between the two environments is the line 
% spacing between each entry. Both environments take one argument, which is the
% title of the section. See main document for how these environments are used.

% Define settings for left-hand column in which dates are printed
\newcolumntype{R}{>{\raggedleft}p{1.3in}}

% Define 'SectionTable' environment
\newenvironment{SectionTable}[1]{
	\renewcommand*{\arraystretch}{1.3} % CHANGE THIS NUMBER FOR TEXT SPACING
	\setlength{\tabcolsep}{10pt}
	\begin{longtable}{Rp{4.9in}} & #1 \\}
{\end{longtable}}

% Define 'SectionTableSingleSpace' environment
\newenvironment{SectionTableSingleSpace}[1]{
	\renewcommand*{\arraystretch}{1.0}
	\setlength{\tabcolsep}{10pt}
	\begin{longtable}{Rp{5.2in}} & #1 \\[0.6em]}
{\end{longtable}\vspace{-.3cm}}

% --- Document starts here ---

\begin{document}
\thispagestyle{firstpage} % This must be right after \begin{document}

% \vspace{0.3cm} % Adjust the vertical spacing before the line
% \hrule height 0.01pt % This adds a horizontal line
\vspace{0.7cm} % Adjust spacing after the line


% --- Name and contact information ---
% Name and date of last update to this document

\noindent{\LARGE{\color{CustomBlue}{Dr. Alexander Nettekoven}}
\hfill{\it\footnotesize Updated \today}}
% \hfill{\it\footnotesize Updated January 16, 2025}}

% % --- Contact information and other items ---

\vspace{0.3cm} 
\begin{center}
% Ignore errors because they don't understand the \ifprivate
\ifprivate{\begin{tabular}{p{2.1in}p{2.1in}p{2.1in}}}{\begin{tabular}{p{1.9in}p{1.7in}p{2.5in}}}
% Line 1: Email, GitHub, office location
\textbf{Role}: Postdoc @ UTexas      &
\textbf{Website}: \href{https://www.alexkoven.com/}{alexkoven.com}    &
\textbf{Email}: \ifprivate{nettekoven@utexas.edu}{alexkoven.inquiry@gmail.com} \\

% Line 2: Phone number, LinkedIn, citizenship
\textbf{US Work}: No Restrictions & 
\textbf{GitHub}: \href{https://github.com/alexkoven}{alexkoven}  &
\textbf{LinkedIn}: \href{https://linkedin.com/in/alexander-nettekoven}{alexander-nettekoven}  
\end{tabular}
\end{center}

\vspace{0.2cm} 

% \vspace{0.3cm} % Adjust the vertical spacing before the line
\hrule height 0.5pt % This adds a horizontal line
% \vspace{0.3cm} % Adjust spacing after the line

% % --- Section: Research interests ---

% \begin{SectionTable}{\headingfont Research Interests}
% & Multi-agent AI systems, Autonomous drones, Robotics, Decision-making AI, Additive manufacturing, Control theory
% \end{SectionTable}


% --- Section: Education ---

\begin{SectionTable}{\headingfont Education}
2019 -- 2022 & 
\textbf{University of Texas at Austin}, Austin, TX \newline
Ph.D. in Mechanical Engineering (Specialization: Robotics) \newline
NSF CPS Researcher (2019--2022) $|$ Advisor: Prof. Ufuk Topcu \\

2016 -- 2018 & 
\textbf{University of Texas at Austin}, Austin, TX \newline
M.Sc. in Mechanical Engineering \\ 

2011 -- 2016 & 
\textbf{RWTH Aachen University}, Aachen, Germany \newline
B.Sc. in Mechanical Engineering \\

2014 -- 2015 &
\textbf{University of Michigan}, Ann Arbor, MI \newline
Exchange Research Scholar, Barton Research Group \\

\end{SectionTable}

\vspace{-0.8cm}

% --- Section: Professional Experience ---

\begin{SectionTable}{\headingfont Professional Experience}

06/2025 -- Present &
\textbf{Center for Autonomy}, UT Austin -- \textbf{Postdoctoral Fellow} \newline
Systematically deploying and evaluating state-of-the-art robot foundation models, including vision-language-action (VLA) and world action models (WAM):
\begin{itemize}[noitemsep, topsep=2pt]
\item Built low-cost ($<$\$1,200 including onboard compute) bimanual manipulator testbed designed for rapid imitation learning iterations; enabling efficient policy learning with world model verifiers
\item Finetuned, deployed, and evaluated DreamZero, \(\pi\)0.5, SmolVLA, ACT on pick-and-place and household cleaning tasks across simulation and real-world settings
\end{itemize} \\

09/2022 -- 06/2025 &
\textbf{Multi AI, Inc.}, Austin, TX / San Francisco, CA -- \textbf{Co-Founder \& CEO} \newline
Built and scaled a multi-agent reinforcement learning platform for fleet coordination problems ($>$100 vehicles) in the US Air Force and Pentagon:
\begin{itemize}[noitemsep, topsep=2pt]
\item Outperformed baseline decision-making by 38\% on key resource delivery metrics in contested logistics; computed in $<$15 min vs.\ $>$24h for human planners
\item Received numerous Air Force SBIR Phase I/II/III contracts and a TACFI selection (7-figure contract funding total)
\item Tech stack: multi-agent RL for fleet coordination and motion planning; LLMs for human-machine interfacing; imitation learning for warm-starting policies with expert demonstrations
\end{itemize} \\

11/2021 -- 09/2022 &
\textbf{Droneconia LLC} (later Multi AI, Inc.), Austin, TX -- \textbf{Co-Founder \& CTO} \newline
Advanced a PhD-stage prototype to secure the first SBIR Phase I contract with Air Force Agility Prime:
\begin{itemize}[noitemsep, topsep=2pt]
\item Developed a rapid multi-agent motion and mission planner with decentralized Monte Carlo Tree Search (MCTS) for autonomous UAV/UGV planning; scaled to 100+ vehicles
\end{itemize} \\

01/2019 -- 12/2022 &
\textbf{Autonomous Systems Group}, UT Austin -- \textbf{NSF-Supported Researcher} \newline
Developed novel robotic systems and algorithms via robot learning, formal methods, and controls in collaboration with Caltech, UC Berkeley, Northeastern, and RPI:
\begin{itemize}[noitemsep, topsep=2pt]
\item Designed and deployed novel hardware testbeds, including the world's first aerial 3D printing hexacopter built with custom MPC and PID/LQR feedback control for in-flight material deposition
\item Achieved 5.43\,mm RMSE positioning accuracy, competitive with ground-based robotic systems
\item Supervised 20+ undergraduate researchers; published at IROS, ICUAS, IEEE CDC, ICMLA, and SFF
\end{itemize} \\

08/2016 -- 12/2018 &
\textbf{Advanced Manufacturing Center}, UT Austin -- \textbf{Research Assistant} \newline
Developed data-driven control (Predictive Iterative Learning Control) and computer vision pipelines for advanced manufacturing processes and robotics applications. \\

01/2016 -- 04/2016 &
\textbf{Mercedes-Benz}, Stuttgart, Germany -- \textbf{Robot Operations Intern} \newline
Optimized ABB robotic assembly for the OM-651 diesel engine at Daimler headquarters; implemented trajectory generation and optimization for 6-DOF robot arms. \\

11/2015 -- 12/2015 &
\textbf{Ford Germany}, Cologne, Germany -- \textbf{Maintenance Intern} \newline
Diagnosed and repaired ABB industrial robots in aluminum die casting production. \\

04/2015 -- 11/2015 &
\textbf{Fraunhofer Institute of Laser Technology}, Aachen, Germany -- \textbf{Research Assistant} \newline
Fabricated iron-aluminide alloy specimens via Laser Metal Deposition (LMD), achieving densities $>$99.5\%; programmed end-effector trajectories and managed laser safety protocols. \\

10/2014 -- 02/2015 &
\textbf{University of Michigan}, Ann Arbor, MI -- \textbf{Undergraduate Research Scholar} \newline
Developed nano-scale signal detection algorithms for Electrohydrodynamic Jet (e-jet) printing using embedded sensing and signal detection devices in Prof. Kira Barton's lab. \\

% 07/2014 -- 10/2014 & 
% \textbf{Machine Tools and Production Engineering Lab} (Aachen, Germany) -- \textbf{Research Assistant};
% Assisting strategic positioning of global production sites.

\end{SectionTable}
\vspace{-0.5cm}

% --- Section: Achievements ---

\begin{SectionTable}{\headingfont Achievements \& Recognitions}



% 2025 & \textbf{AFLCMC/WLQ and Force Design, Integration, and Wargaming Directorate} \\

2024 & \textbf{Pentagon Wargaming Development Partnership}. Six-figure SBIR Phase III by Chief of Wargaming, OUSD A-S (Pentagon), and Air Force Autonomy Prime. \\

2024 & \textbf{Invited Expert to Texas AI Caucus}. Advised Texas State Representative Giovanni Capriglione on AI policy and commercial regulations. \\

2024 & \textbf{Seven-Figure TACFI Selection by US Air Force}. Awarded for highly-selective Funding Increase program by Air Force Ventures for decision-support AI. \\

2023 & \textbf{Major Phase III Award by US Air Force Autonomy Prime}. Six-figure contract for multi-agent coordination AI. \\

% 2023 & \textbf{Expert Speaker on Generative AI}. Industry panelist at Applied Intelligence Summit 2023. \\

2023 & \textbf{Capital Factory Accelerator}. Selected for Texas's largest and most active startup accelerator. \\

2022 & \textbf{Launch Texas Awardee}. Launch Texas awardee of newly-founded space entrepreneurship program at UT Austin. \\

2022 & \textbf{Major Phase II Award by US Air Force Agility Prime}. Secured first six-figure contract for multi-agent coordination AI to accelerate rapid scaling in DoD. \\

2022 & \textbf{Selection for Elite Startup Accelerators}. MassChallenge finalist out of 400 startups and selection to exclusive Air Force Technology Acceleration Program. \\

2022 & \textbf{Texas Venture Labs Commercialization Award Winner}. Five-figure prize for highest commercialization success among Texas student startups. \\

2022 & \textbf{Bridge-to-Phase II Award Winner}. UT Discovery grant for accelerating national security decision-support AI. \\

2022 & \textbf{Major Air Force Programs and Prime Support}. Formal endorsements from AFWERX, AFRL ACT3, RQQ, RISB, and Textron for startup's solution offering. \\

2021 & \textbf{GAIN Mechanical Engineering Department Award}. Outstanding research contributions award at University of Texas at Austin. \\

2019--2022 & \textbf{National Science Foundation (NSF)-Supported Researcher}. Selected researcher on multi-year project with CalTech, Berkeley, Northeastern, and RPI. \\

2017 & \textbf{Sandia National Laboratories Grant}. Proposed and won five-figure funding for innovative laser pre-scan process enhancing 3D printing quality. \\

2015 & \textbf{Undergraduate Research Opportunity Award}. Dual award by the University of Michigan and RWTH Aachen University for undergraduate research. \\

\end{SectionTable}

\vspace{-0.5cm}

% --- Section: Skills ---

\begin{SectionTable}{\headingfont Skills}

& \textbf{Robotics \& Control:} Manipulation, Robot Deployment, Reinforcement Learning (PPO), Deep Learning, Perception (SLAM, VIO), Motion Planning (trajectory generation, trajectory optimization), State Estimation (Kalman Filters, EKF, EKF2), Feedback Control (PID, Optimal Control, Nonlinear Control), Model Predictive Control (MPC) \\

& \textbf{VLA \& Foundation Models:} World Models, \(\pi\)0.5, SmolVLA, DreamZero, ACT, Diffusion Policies, Imitation Learning (DAgger, teacher-student learning), Sim-to-Real Transfer \\

& \textbf{Programming \& Tools:} Python, C++, PyTorch, ROS2, Isaac Sim/Isaac Lab, Docker, Linux, Edge Deployment, Systems Design, Git \\

& \textbf{Team:} Product and project management; technical leadership; team building and hiring; technical communication and presentations \\

& \textbf{Languages:} German (native), English (native) \\

\end{SectionTable}

\vspace{-0.5cm}

% --- Section: Publications ---

\begin{SectionTable}{\headingfont Publications}

2024 &
\textbf{Alexander Nettekoven}, Matthew Carlin, and Ufuk Topcu. ``Agile Combat Employment System (ACES) Model Modifications.'' \textit{Technical Report to US Air Force AFWERX SBIR Phase III}. \\

2023 &
\textbf{Alexander Nettekoven}, Matthew Carlin, and Ufuk Topcu. ``Large-Scale Fleet Coordination for Contested Logistics.'' \textit{Technical Report to US Air Force AFWERX SBIR Phase II}. \\

2022 &
\textbf{Alexander Nettekoven}. ``On the Feasibility of 3D Printing with Multicopters.'' \textit{Ph.D. Dissertation at the University of Texas at Austin}. \\

2022 &
\textbf{Alexander Nettekoven}, Scott Fish, Joseph Beaman, and Ufuk Topcu. ``Towards online monitoring and data-driven control: a study of segmentation algorithms for infrared images of the powder bed.'' \textit{2022 International Solid Freeform Fabrication Symposium}. \\

2021 &
\textbf{Alexander Nettekoven}, Nicholas Franken, and Ufuk Topcu. ``Geometrical Analysis of Simple Contours Deposited by a 3D Printing Hexacopter.'' \textit{2021 International Solid Freeform Fabrication Symposium}. \\

2021 &
\textbf{Alexander Nettekoven} and Ufuk Topcu. ``A 3D Printing Hexacopter: Design and Documentation.'' \textit{2021 International Conference on Unmanned Aircraft Systems (ICUAS)}. \\

2021 &
Zizhao Wang, Xuesu Xiao, \textbf{Alexander J. Nettekoven}, Kadhiravan Umasankar, Anika Singh, Sriram Bommakanti, Ufuk Topcu, and Peter Stone. ``From agile ground to aerial navigation: Learning from learned hallucination.'' \textit{2021 IEEE/RSJ International Conference on Intelligent Robots and Systems (IROS)}. \\

2019 &
Yingshui Tan, Baihong Jin, \textbf{Alexander Nettekoven}, Yuxin Chen, Yisong Yue, Ufuk Topcu, and Alberto Sangiovanni-Vincentelli. ``An encoder-decoder based approach for anomaly detection with application in additive manufacturing.'' \textit{2019 18th IEEE International Conference on Machine Learning and Applications (ICMLA)}. \\

2019 &
Zhe Xu, \textbf{Alexander J. Nettekoven}, A. Agung Julius, and Ufuk Topcu. ``Graph temporal logic inference for classification and identification.'' \textit{2019 IEEE 58th Conference on Decision and Control (CDC)}. \\

2018 &
\textbf{Alexander Nettekoven}. ``Predictive iterative learning control with data-driven model for near-optimal laser power in selective laser sintering.'' \textit{2018 Master's Thesis}. \\

2018 &
\textbf{Alexander Nettekoven}, Scott Fish, Joseph Beaman, and Ufuk Topcu. ``Predictive Iterative Learning Control with Data-Driven Model for Optimal Laser Power in Selective Laser Sintering.'' \textit{2018 International Solid Freeform Fabrication Symposium}.
\end{SectionTable}

\vspace{-0.5cm}


% --- Section: Services ---

\begin{SectionTable}{\headingfont Services}
2021 -- present & \textbf{Peer Reviewer for Robotics and Engineering Conferences:}
\begin{itemize}[noitemsep, topsep=0pt]
\item International Conference on Robotics and Automation (ICRA)
\item International Conference on Intelligent Robots and Systems (IROS)
\item IEEE Robotics and Automation Letters (RA-L)
\item American Control Conference (ACC)
\item North American Manufacturing Research Conference (NAMRC)
\item Solid Freeform Fabrication Symposium (SFF)
\item Conference on Automation Science and Engineering (CASE)
\end{itemize} \\

2020 -- present & \textbf{Guest Lecturer for Undergraduate Feedback Control}. Presentations of real-world control applications to engineering students at UT Austin. \\

2025 & \textbf{STEM Outreach, Del Valle ISD}. Taught students to build and program a self-driving RC car. \\

2021 & \textbf{Texas Robotics Motion Capture Space}. Designed, directed, and commissioned high-resolution motion capture space for UT Austin robotics program. \\

2021 & \textbf{Public Safety Drone Collaboration Workshop}. Organized Drone workshop for Austin PD, Texas DPS, Travis County Sheriff's Office, and Austin Fire. \\

2020 -- 2022 & \textbf{Research Mentor}. Supervised 7 undergraduate research assistants: Sriram Bommakanti, Kadhir Umasankar, Nicholas Franken, Harrison Jin, Meenakshi Swaminathan, Steven Diep, and Mario Gonzalez. \\

2020 & \textbf{Research Instructor}. Led full-time summer research internships for 17 undergraduate students in the Autonomous Systems Group. \\

\end{SectionTable}

\vspace{-0.5cm}

\begin{SectionTable}{\headingfont Press Coverage}

2023 & Jeremey R. Dunn. "Tech Connect partner uses AI to assist in in-theater mission planning." \textit{Defense Visual Information Distribution Service} (08/23/2023). \\

2023 & Deborah Yao. "Hybrid AI Reduces Generative AI Hallucinations, Applied Intelligence Live! Austin 2023." \textit{AI Business} (09/21/2023). \\

2022 & Nat Levy. "Ready for Launch." \textit{Texas Engineer} (2022-2023 issue).

\end{SectionTable}

\vspace{-1.0cm}


\begin{SectionTable}{\headingfont Invited Talks and Presentations}
12/2024 & US Air Force Wargaming Offices and OUSD A-S Chief of Wargaming Workshop. \\

03/2024 & HAF/a5, SAF/SA, HAF/A4, RGPA, AMC A5-8, AMC A9 Wargaming Capabiltiy Review. \\

09/2023 & Capital Factory Accelerator First Look Kickoff. \\

08/2023 & National Security Innovation Council (NSIC). \\

05/2023 & Texas Innovation Center Board Meeting and Luncheon. \\

11/2022 & Launch Texas Award and Kickoff at the University of Texas at Austin. \\

09/2022 & Blue Tech Collider Multi-Drone Show Case. \\

08/2022 & Solid Freeform Fabrication Conference 2022. \\

07/2022 & MassChallenge Early-Stage Startup Accelerator Kickoff. \\

02/2022 & Hanscom and Incirlik US Air Force Base Capability Review. \\

08/2021 & Solid Freeform Fabrication Conference 2021. \\

06/2021 & International Conference on Unmanned Aircraft Systems. \\

04/2021 & Texas Robotics Research Seminar. \\

02/2021 & Graduate and Industry Networking event, University of Texas at Austin. \\

10/2020 & Walker Department of Mechanical Engineering, University of Texas at Austin. \\

03/2019 & NSF Cyberphysical System Review with Caltech, Berkeley, Northeastern, RPI. \\

08/2018 & Solid Freeform Fabrication Conference 2018. \\

02/2017 & America Makes Project Review with US Air Force AFRL. \\

06/2016 & International Center for Sustainable Textiles at RWTH Aachen University. \\

02/2015 & Barton Research Group at the University of Michigan. \\

\end{SectionTable}


% --- End of CV! ---

\end{document}






